%% =====================================================================
%%  paper2_introduction_v3.tex
%%
%%  Revised Introduction (Section 1) -- v3.
%%  Drop-in replacement for the \section{Introduction} block at lines
%%  120--218 of paper2_delay_cascade.tex.  Compiles against cas-dc.cls.
%%
%%  ---- What changed from v2 (and WHY, from a reviewer's standpoint) ----
%%
%%  R1. Removed in-Introduction \subsection*{...} headings.  Elsevier
%%      EPSR / RESS / Applied Energy reviewers expect a single flowing
%%      Introduction; sub-headings inside Section 1 read as a thesis
%%      chapter rather than a journal paper.
%%
%%  R2. Removed the four-row T0/T1/T2/T3 comparison TABLE and the
%%      label "T3 = This work" inside it.  Putting our own work in a
%%      taxonomy table is interpreted as self-promotional.  The same
%%      taxonomy is now narrated in two prose paragraphs and our gap
%%      is stated as a logical consequence, not a line in our own table.
%%
%%  R3. Cut the contribution list from 5 items to 4.  Conventional EPSR
%%      Introductions list 3-4 contributions; 5 invariably triggers the
%%      reviewer comment "the contributions are not focused".  The UFLS
%%      over-shedding mechanism is now folded into the closed-loop
%%      contribution as a sub-bullet, not a separate top-level claim.
%%
%%  R4. Removed the inline \begin{quote}...\end{quote} feedback-loop
%%      block and replaced it with a single formal sentence.  Block
%%      quotes inside Introductions of engineering journals are rare
%%      and look essayistic.
%%
%%  R5. Reduced the standard-anchor citation density.  v2 cited six
%%      named standards in one paragraph (PRC-002-2, PRC-006-5,
%%      CIP-012, NASPI, C37.118.2, IEC 61850 PRP), which reads as
%%      name-dropping.  v3 keeps three central anchors in the
%%      Introduction (NERC PRC-002-2, NASPI, IEC 61850 PRP) and defers
%%      the others to the Methods section where they are actually used.
%%
%%  R6. Added a single explicit research question at the end of the
%%      gap paragraph ("RQ"), so a reviewer immediately knows the
%%      claim that the paper will be evaluated against.
%%
%%  R7. Softened or removed the unsupported claim that T2 misses
%%      "interaction between channels"; the gap is now stated only in
%%      terms of what we DO model, not what others fail to model.
%%
%%  R8. Removed Xin2024CPPSframework and Wang2022CPPS until their
%%      bibliographic metadata can be hand-verified.  v3 cites only
%%      references whose volume/issue/year are confidently correct.
%%      The T2 description is now anchored on Cai2016TSG and
%%      LaiCPPS2019, which are already in the original bib.
%%
%%  R9. Highlights and Abstract elsewhere in the .tex file should be
%%      synchronised with the four contributions listed below.  See the
%%      end-of-file comment block for the recommended re-wording.
%% =====================================================================

\section{Introduction}\label{sec:intro}

The deep integration of physical infrastructure with high-speed
communication networks has transformed traditional power grids into
\emph{cyber-physical power systems} (CPPS), in which a physical power
layer and a cyber communication layer operate as a tightly coupled
whole \citep{Rinaldi2001,Ouyang2014}.  Under nominal conditions this
coupling improves dispatch flexibility, situational awareness and
wide-area control performance.  Under stress, however, it opens a
bidirectional failure channel: a disturbance originating in either
layer can propagate to the other and trigger cascading failures of a
magnitude neither layer would experience in isolation
\citep{Buldyrev2010,Gao2012}.

A substantial body of work has investigated such failures through
topological percolation models, flow-redistribution rules and
attack-vulnerability analyses
\citep{Motter2002,Dobson2007,Brummitt2012,Schneider2011}.  These
studies established \emph{how} structural interdependence amplifies
failure propagation, and they remain the conceptual backbone of
modern CPPS robustness analysis.  In most of them, however, the
communication layer is represented as a binary entity---a link is
either functional or failed---without modelling the quality of
service that surviving links provide.  When end-to-end latency is
addressed at all, it is either absent from the cascade dynamics or
applied as a post-hoc rescaling of the final performance metric, so
that all delay scenarios share an identical structural cascade
\citep{Milano2012,LaiCPPS2019}.  More recent work has begun to
relax this binary view by modelling a single delay channel
explicitly---typically queueing on transmission links or jitter on
PMU streams---and embedding it inside the cascade
\citep{Cai2016TSG,Falahati2014TSG,Huang2015}.  This direction is the
right one, but the chain that feeds modern wide-area protection is
not single-channel.  Industrial measurement budgets distinguish
sensor and PMU acquisition, edge processing, transmission and
actuation as separately bounded components
\citep{NASPI_LatencyGuidelines,IEEE_C37_118_2}, and the redundancy
practices that sit on top of them
\citep{IEC_61850_PRP} suppress the loss probability of each
component to very different degrees.  A model in which all of these
components collapse into a scalar~$\tau$ therefore cannot, by
construction, reproduce the asymmetric way real cyber stress maps
into closed-loop control degradation.

This observation motivates the central research question of the
present paper:
\begin{quote}
\emph{Does embedding a standard-anchored, multi-component
communication delay model inside every round of a CPPS cascade
produce qualitatively different vulnerability conclusions from a
delay-free or single-$\tau$ model on the same network?}
\end{quote}

We answer this question affirmatively, on three counts that we
develop in the body of the paper: the delay penalty is asymmetrically
concentrated at moderate-to-high tolerance margins (precisely the
operating regime in which a delay-free analysis predicts safety);
delay amplifies the stochastic variability of cascade outcomes,
producing heavy-tailed terminal distributions that ideal-communication
models do not exhibit; and a generator-fidelity metric exposes
functional degradation that aggregate load-retention metrics fail to
detect.  The mechanism that produces all three effects is a
self-reinforcing loop: degraded communication reduces the closed-loop
control efficacy of every surviving generator, the resulting
imbalance triggers additional branch overloads, the structural
failures further damage the cyber topology, and the damaged topology
amplifies the latency that drives the next round.  Our prior work
\citep{ChenTu2026} established the underlying CPPS framework with
distributed control centers, a Functional Viability-Constrained
Island Detection mechanism and a Cohesion-Propagation Importance
Metric for control-center placement, but did not model the
influence of communication delay on generator control during the
cascade itself.  The present paper closes that gap.

The contributions of this work are as follows.

\begin{enumerate}

\item \textbf{A delay-integrated, closed-loop cascading failure
  framework.}  Communication delay is computed inside every cascade
  round and applied to generator output \emph{before} the
  DC power-flow redistribution that determines branch overloads.
  Each delay scenario therefore evolves along a genuinely
  different trajectory---different overload patterns, different
  failure sequences, different terminal states---rather than along
  a rescaled copy of a common delay-free trajectory.  The same
  loop also drives an underfrequency load-shedding (UFLS) rule
  whose shed depth is a function of the prevailing latency, so
  that the standard NERC PRC-006-5 single-stage rule is recovered
  as the zero-delay limit.

\item \textbf{A four-component, role-aware end-to-end delay
  decomposition.}  Latency is split into measurement
  ($\tau_m$), edge processing ($\tau_e$), transmission ($\tau_t$)
  and actuation ($\tau_a$) components, with role-dependent service
  times for control centers and ordinary relay nodes, mirroring
  the asymmetric processing behaviour of SCADA / EMS architectures.
  Each component is parameterised against an externally defensible
  reference: NASPI synchrophasor latency
  guidelines~\citep{NASPI_LatencyGuidelines} and
  IEEE~C37.118.2~\citep{IEEE_C37_118_2} bound $\tau_m$ and
  $\tau_t$, the IEC~61850 PRP / HSR redundancy
  stack~\citep{IEC_61850_PRP} bounds the residual loss
  probability, and the NERC PRC-002-2 sub-second budget
  \citep{NERC_PRC002} bounds the heaviest scenario we test.

\item \textbf{A three-mechanism control-efficacy function
  $\boldsymbol{\eta^{+} = \Phi_{\mathrm{loss}} \cdot
  \Phi_{\mathrm{sat}} \cdot \Phi_{\mathrm{crit}}}$.}  Delay-induced
  generator-control loss is decomposed into three orthogonal
  mechanisms with separately tunable parameters:
  redundancy-attenuated message loss
  ($\Phi_{\mathrm{loss}}$), continuous controller saturation under
  elevated $\tau_m+\tau_e$ ($\Phi_{\mathrm{sat}}$) and a normalised
  collapse near each machine's critical
  delay~$\tau_{\mathrm{crit},i}$ ($\Phi_{\mathrm{crit}}$, normalised
  so that $\eta^{+}(\tau{=}0) = 1$ exactly).  Because the three
  mechanisms are orthogonal in parameter space, any observed
  vulnerability can be attributed to a specific physical channel,
  rather than absorbed into a single black-box delay penalty.

\item \textbf{Twin resilience metrics that reveal the asymmetric
  delay-penalty regime and its hidden component.}  The cascade
  is evaluated jointly on a supply-conservation ratio~$R_1$ and on
  a capacity-weighted dispatch-fidelity NRMSE~$R_3$.  A
  delay-penalty heatmap on the (tolerance, cascade-round) plane
  shows that the damage is largest at moderate-to-high tolerance
  margins, narrowing the critical intervention window to the
  early cascade rounds, while $R_3$ exposes generator-level
  functional degradation that $R_1$ alone misses.  These results
  show that delay-free and single-$\tau$ robustness models are
  not merely imprecise but systematically misleading in the
  operating regime that operators most often rely on.

\end{enumerate}

The remainder of this paper is organised as follows.
Section~\ref{sec:model} presents the CPPS model, including the
four-component communication delay formulation and the
betweenness-assortative interdependent coupling.
Section~\ref{sec:cascade} describes the delay-integrated cascading
failure dynamics, the $\eta^{+}$ control-efficacy function, the
delay-coupled UFLS rule and the five delay-severity scenarios.
Section~\ref{sec:experiment} reports Monte~Carlo experiments on the
IEEE~39-bus test system coupled to a Barab\'asi--Albert scale-free
communication network, covering 1\,000 coupling realisations,
eleven tolerance levels and the five delay scenarios, and analyses
the asymmetric delay-penalty regime, the heavy-tailed collapse
distributions, the hidden $R_3$ degradation, the round-level
intervention window and a sensitivity analysis of mitigation
actions.  Section~\ref{sec:conclusion} concludes the paper.

%% ---------------------------------------------------------------------
%%  Recommended downstream edits to keep the front-matter consistent
%%  with the four-contribution structure above:
%%
%%  * Highlights bullet 5 ("Round-level heatmap analysis...") should
%%    be merged into bullet 4 to avoid foregrounding a result before
%%    the contribution list does.
%%  * The Abstract sentence "delay decomposed into uplink ... and
%%    downlink ..." should mention the four named components
%%    (measurement, edge, transmission, actuation) so the abstract
%%    matches Contribution 2.
%% ---------------------------------------------------------------------
